% =============================================================
%  07_backtracking.tex
%  Capitolo 7 -- Backtracking e algoritmi
% =============================================================

\chapter{Backtracking e algoritmi}

\begin{intuizione}
Nei capitoli precedenti le funzioni trasformavano dati:
una lista entra, una lista diversa esce.
Qui le funzioni \emph{esplorano possibilità}: provano una scelta,
verificano se va bene, e se non va tornano indietro
per provarne un'altra.

\medskip
Questo schema si chiama \textbf{backtracking}.
Non è un algoritmo specifico: è un modo di strutturare
la ricerca in uno spazio di soluzioni.
La ricorsione che hai imparato nel Capitolo~3 è lo strumento naturale
per esprimerlo.
\end{intuizione}

\begin{tcolorbox}[
  enhanced,
  colback=black!2, colframe=black!20,
  arc=3pt, boxrule=0.4pt,
  left=12pt, right=12pt, top=9pt, bottom=9pt,
  before skip=8pt, after skip=10pt
]
\small\color{black!65}

\textbf{Il percorso di questo capitolo.}

\medskip
\begin{enumerate}
  \item \textbf{Generate-and-test.} Il modello più semplice:
        genera tutti i candidati, testa quelli che soddisfano il vincolo.
  \item \textbf{Stato, scelta, vincolo.} I tre ingredienti di
        ogni algoritmo di backtracking.
  \item \textbf{Pruning.} Tagliare rami prima di esplorarli:
        trasforma un algoritmo esponenziale in qualcosa di accettabile.
  \item \textbf{Permutazioni e sottoinsiemi.} Due problemi classici
        come istanza dello stesso schema.
  \item \textbf{N-Queens.} Il problema archetipico del backtracking,
        costruito passo per passo.
  \item \textbf{Sudoku.} Lo stesso modello su un dominio più ricco.
\end{enumerate}
\end{tcolorbox}

\vspace{4pt}

% ================================================================
\section{Generate-and-test}
% ================================================================

Il modo più diretto per trovare soluzioni in uno spazio discreto
è questo: genera tutti i candidati possibili, poi tieni solo
quelli che soddisfano il vincolo.
Non è efficiente, ma è corretto e facile da ragionare.

\begin{concetto}{Generate-and-test}
\keyline{Genera}{produci tutti i candidati (o un sottoinsieme).}
\keyline{Testa}{per ognuno, verifica il predicato di correttezza.}
\keyline{Raccogli}{tieni i candidati che passano il test.}

\boxsep
\detail{È il pattern \fn{remove-if-not} applicato a uno spazio
di soluzioni invece che a una lista già pronta.
La difficoltà sta nel definire correttamente cosa si
intende per ``tutti i candidati''.}
\end{concetto}

\begin{esempio}[Trova le coppie la cui somma è N]
\begin{lstlisting}
;; Genera tutte le coppie (i j) con i < j da una lista,
;; poi filtra quelle che sommano a target.

(defun coppie-con-somma (lst target)
  (let ((coppie '()))
    (dolist (i lst)
      (dolist (j lst)
        (when (and (< i j) (= (+ i j) target))
          (push (list i j) coppie))))
    (nreverse coppie)))

(coppie-con-somma '(1 2 3 4 5 6) 7)
; => ((1 6) (2 5) (3 4))
\end{lstlisting}
\end{esempio}

\begin{nota}
Generate-and-test con loop esplicito va bene per spazi piccoli.
Quando lo spazio cresce in modo combinatorio (es. tutti i sottoinsiemi
di una lista di 20 elementi sono $2^{20} \approx 10^6$), serve
qualcosa di più intelligente: il backtracking vero.
\end{nota}

\begin{workbook}
\textbf{Esercizi 121--122}\\[4pt]
L'esercizio~121 (generate-and-test su interi) è il punto di entrata:
prendi confidenza con il pattern prima di aggiungere la ricorsione.
\end{workbook}

% ================================================================
\section{Stato, scelta, vincolo}
% ================================================================

Ogni algoritmo di backtracking ha tre ingredienti.
Riconoscerli nel problema è il primo passo per implementarlo.

\argomento{Stato}

\begin{concetto}{Stato}
\keyline{}{Lo stato è la configurazione parziale costruita fin qui.}
\keyline{}{All'inizio è vuoto. Ad ogni passo cresce di una scelta.}

\boxsep
\detail{In Lisp lo stato è tipicamente una lista delle scelte
fatte finora, passata come parametro ricorsivo.
Non serve modificarla: si costruisce una nuova lista
ad ogni chiamata ricorsiva.}
\end{concetto}

\begin{esempio}[Stato come lista di scelte]
\begin{lstlisting}
;; Genera tutte le sequenze binarie di lunghezza n.
;; Lo "stato" e' la lista dei bit scelti finora.

(defun binarie (n &optional (stato '()))
  (if (= (length stato) n)
      (list stato)                        ; soluzione completa: raccoglila
      (append
        (binarie n (cons 0 stato))        ; scelta: aggiungi 0
        (binarie n (cons 1 stato)))))     ; scelta: aggiungi 1

(binarie 2)
; => ((0 0) (0 1) (1 0) (1 1))
\end{lstlisting}
\end{esempio}

\argomento{Scelta}

\begin{concetto}{Scelta}
\keyline{}{Ad ogni passo si sceglie un elemento da un insieme di candidati.}
\keyline{}{La funzione ricorsiva viene chiamata una volta per ogni scelta.}

\boxsep
\detail{L'insieme dei candidati dipende dallo stato corrente:
scelte già fatte possono escludere certi valori.
Questo è il punto in cui si inserisce il pruning.}
\end{concetto}

\argomento{Vincolo}

\begin{concetto}{Vincolo}
\keyline{}{Il vincolo dice quando uno stato è una soluzione valida.}
\keyline{}{Può essere controllato alla fine (stato completo) o durante (pruning).}
\end{concetto}

\begin{intuizione}
Un algoritmo di backtracking è una funzione ricorsiva
che esplora un albero di scelte.
Ogni nodo è uno stato. Ogni arco è una scelta.
Le foglie sono gli stati completi. Il vincolo seleziona
le foglie valide.

\medskip
Se il vincolo viene verificato \emph{anche nei nodi interni},
si tagliano interi sottoalberi prima di raggiungerli.
Questo è il pruning.
\end{intuizione}

\begin{confronto}
\textbf{Confronto con il backtracking imperativo (Java/C).}
In un linguaggio imperativo il backtracking si implementa
tipicamente con cicli annidati e \texttt{continue}/\texttt{break},
o con una struttura di stack esplicita.
La versione ricorsiva in Lisp usa lo stack delle chiamate
come stack implicito: il ``tornare indietro'' è semplicemente
il ritorno dalla chiamata ricorsiva.
Non serve stato globale mutabile: ogni ramo dell'albero
ha il proprio stato locale.
\end{confronto}

\begin{workbook}
\textbf{Esercizi 123--124}\\[4pt]
L'esercizio~123 chiede di enumerare le sequenze con un vincolo
semplice. È il punto in cui padroneggiare la struttura
stato/scelta/vincolo prima di aggiungere il pruning.
\end{workbook}

% ================================================================
\section{Pruning}
% ================================================================

Il backtracking puro esplora tutto lo spazio.
Per molti problemi questo è troppo lento.
Il pruning taglia i rami che non possono portare a soluzioni.

\begin{concetto}{Pruning}
\keyline{}{Prima di fare una scelta, verifica se lo stato parziale
può ancora portare a una soluzione.}
\keyline{}{Se non può, non espandere quel ramo: ritorna \fn{nil}.}

\boxsep
\detail{Il pruning non cambia le soluzioni trovate:
cambia solo il tempo per trovarle.
Un vincolo applicato presto (nodo interno) è sempre
meglio dello stesso vincolo applicato tardi (foglia).}
\end{concetto}

\begin{esempio}[Sequenze senza due 1 consecutivi, con pruning]
\begin{lstlisting}
;; Genera sequenze binarie di lunghezza n
;; senza due 1 consecutivi.
;;
;; Pruning: se l'ultimo elemento aggiunto e' 1,
;; non aggiungere un altro 1 (lo scarto prima della ricorsione).

(defun binarie-no-11 (n &optional (stato '()))
  (if (= (length stato) n)
      (list stato)
      (let ((candidati
              (if (and stato (= (first stato) 1))
                  '(0)          ; l'ultimo e' 1: posso solo aggiungere 0
                  '(0 1))))     ; altrimenti entrambi
        (mapcan (lambda (bit)
                  (binarie-no-11 n (cons bit stato)))
                candidati))))

(binarie-no-11 3)
; => ((0 0 0) (0 0 1) (0 1 0) (1 0 0) (1 0 1))
\end{lstlisting}
\end{esempio}

\begin{attenzione}
Il pruning cambia l'ordine in cui le soluzioni vengono
trovate, ma non quali soluzioni vengono trovate.
Se il pruning è \emph{corretto} (non esclude stati che
potrebbero diventare soluzioni valide), la correttezza
dell'algoritmo è preservata.
Un pruning sbagliato può eliminare soluzioni valide.
Verifica sempre: ``questo vincolo, applicato a uno stato parziale,
esclude davvero tutti gli stati completi non validi derivabili da esso?''
\end{attenzione}

\begin{workbook}
\textbf{Esercizi 124--125}\\[4pt]
L'esercizio~125 richiede di aggiungere pruning a una
funzione di ricerca esistente. Misura il numero di
nodi esplorati prima e dopo: la differenza è spesso
di ordini di grandezza.
\end{workbook}

% ================================================================
\section{Permutazioni e sottoinsiemi}
% ================================================================

Due problemi classici che appaiono frequentemente
e che sono entrambi istanze dirette del modello
stato/scelta/vincolo.

\argomento{Permutazioni}

\begin{concetto}{Permutazioni}
\keyline{}{Una permutazione di \fn{(a b c)} usa ogni elemento esattamente una volta.}
\keyline{}{Stato: elementi già scelti. Scelta: un elemento non ancora usato.}
\end{concetto}

\begin{esempio}[Tutte le permutazioni di una lista]
\begin{lstlisting}
(defun permutazioni (lst)
  (if (null lst)
      '(())                          ; caso base: lista vuota -> una soluzione vuota
      (mapcan
        (lambda (el)
          (let ((resto (remove el lst :count 1)))
            (mapcar (lambda (perm) (cons el perm))
                    (permutazioni resto))))
        lst)))

(permutazioni '(1 2 3))
; => ((1 2 3) (1 3 2) (2 1 3) (2 3 1) (3 1 2) (3 2 1))
\end{lstlisting}
\end{esempio}

\argomento{Sottoinsiemi (power set)}

\begin{concetto}{Sottoinsiemi}
\keyline{}{Il power set di una lista è l'insieme di tutti i suoi sottoinsiemi.}
\keyline{}{Per ogni elemento: due scelte, includerlo o no.}
\keyline{}{Una lista di $n$ elementi ha $2^n$ sottoinsiemi.}
\end{concetto}

\begin{esempio}[Power set]
\begin{lstlisting}
(defun power-set (lst)
  (if (null lst)
      '(())                          ; caso base: il solo sottoinsieme e' nil
      (let* ((testa  (car lst))
             (coda   (cdr lst))
             (senza  (power-set coda))    ; sottoinsiemi senza la testa
             (con    (mapcar (lambda (s) (cons testa s))
                             senza)))    ; sottoinsiemi con la testa
        (append senza con))))

(power-set '(a b c))
; => (NIL (C) (B) (B C) (A) (A C) (A B) (A B C))
\end{lstlisting}
\end{esempio}

\begin{intuizione}
Permutazioni e power set sembrano problemi diversi,
ma hanno la stessa struttura: a ogni passo scegli
un elemento (o se includerlo), e chiami ricorsivamente
su ciò che resta.
La funzione \fn{mapcan} è lo strumento naturale
per raccogliere i risultati di tutte le scelte.
\end{intuizione}

\begin{workbook}
\textbf{Esercizi 126--128}\\[4pt]
L'esercizio~126 (permutazioni) e il~127 (power set) sono
implementazioni dirette di quanto visto qui.
L'esercizio~128 aggiunge un vincolo: genera solo i sottoinsiemi
di taglia $k$. Provalo prima senza guardare la soluzione.
\end{workbook}

% ================================================================
\section{N-Queens}
% ================================================================

Il problema delle N regine è il problema archetipico
del backtracking: posizionare $N$ regine su una scacchiera
$N \times N$ in modo che nessuna si minacci.

\argomento{Il modello}

Prima di scrivere codice, identifica i tre ingredienti.

\begin{concetto}{N-Queens: stato, scelta, vincolo}
\keyline{Stato}{lista delle colonne occupate, una per riga, dall'alto.
Se lo stato è \fn{(3 1)}, la riga~1 ha la regina in colonna~3,
la riga~2 in colonna~1.}
\keyline{Scelta}{la colonna in cui posizionare la regina nella prossima riga.}
\keyline{Vincolo}{nessuna regina minaccia un'altra: stessa colonna
o stessa diagonale.}

\boxsep
\detail{Rappresentare lo stato come lista di colonne (una per riga)
elimina automaticamente il conflitto di riga: ogni riga ha
esattamente una regina.}
\end{concetto}

\argomento{Il predicato di sicurezza}

\begin{esempio}[Verifica se una posizione è sicura]
\begin{lstlisting}
;; Controlla se aggiungere una regina nella colonna `col`
;; (riga corrente = (length stato) + 1) e' sicuro
;; rispetto alle regine gia' piazzate in `stato`.
;;
;; stato: lista delle colonne (prima riga = last, ultima = first).
;; Riga della regina da piazzare: (length stato).

(defun sicura-p (col stato)
  (loop for c in stato
        for dist from 1
        never (or (= c col)               ; stessa colonna
                  (= (abs (- c col))      ; stessa diagonale
                     dist))))
\end{lstlisting}
\end{esempio}

\argomento{La funzione principale}

\begin{esempio}[N-Queens completo]
\begin{lstlisting}
(defun n-queens (n &optional (stato '()))
  (if (= (length stato) n)
      (list stato)                        ; trovata una soluzione
      (mapcan
        (lambda (col)
          (when (sicura-p col stato)      ; pruning: colonna sicura?
            (n-queens n (cons col stato))))
        (loop for i from 1 to n collect i))))

;; Tutte le soluzioni per 4 regine
(n-queens 4)
; => ((3 1 4 2) (2 4 1 3))

;; Solo il numero di soluzioni
(length (n-queens 8))
; => 92
\end{lstlisting}
\end{esempio}

\begin{nota}
Il pruning è dentro il \fn{when}: se \fn{sicura-p} restituisce \fn{nil},
\fn{when} non chiama \fn{n-queens} e \fn{mapcan} raccoglie \fn{nil}
(che viene ignorato). Il taglio avviene senza codice aggiuntivo.

\medskip
La soluzione per $N=8$ trova 92 configurazioni esplorando
una piccola frazione dei $8! = 40320$ candidati possibili.
\end{nota}

\begin{confronto}
\textbf{N-Queens in Java.}
La versione Java tipica usa un array mutabile e un ciclo
\texttt{for} con \texttt{backtrack} esplicito:
si piazza, si verifica, si rimuove.
La versione Lisp non modifica niente: ogni chiamata ricorsiva
crea un nuovo stato (la lista \fn{(cons col stato)}) senza
toccare lo stato precedente.
Il ``backtrack'' non è un'operazione esplicita:
è semplicemente il ritorno dalla chiamata ricorsiva
che non ha trovato soluzioni.
\end{confronto}

\begin{workbook}
\textbf{Esercizi 129--132}\\[4pt]
L'esercizio~129 implementa N-Queens seguendo lo schema di questo
capitolo. Gli esercizi~130--132 lo estendono: conta le soluzioni,
stampa la scacchiera, trova solo la prima soluzione.\\[4pt]
L'esercizio~131 (``prima soluzione'') è particolarmente istruttivo:
mostra come usare le eccezioni per uscire dalla ricerca appena
trovata la prima risposta valida.
\end{workbook}

% ================================================================
\section{Sudoku}
% ================================================================

Il Sudoku è un'istanza più ricca dello stesso modello.
Lo spazio di ricerca è più grande, ma la struttura
dell'algoritmo non cambia.

\begin{concetto}{Sudoku come backtracking}
\keyline{Stato}{la griglia $9 \times 9$ con alcune celle già riempite.}
\keyline{Scelta}{scegli la prima cella vuota; prova i valori 1--9.}
\keyline{Vincolo}{il valore non deve apparire nella stessa riga,
colonna o box $3 \times 3$.}

\boxsep
\detail{In Lisp la griglia si rappresenta comodamente come vettore
di 81 elementi (accessibile con \fn{aref}) o come lista di liste.
Il vettore è più efficiente per l'accesso per indice.}
\end{concetto}

\begin{esempio}[Schema di un risolutore Sudoku]
\begin{lstlisting}
;; Rappresentazione: vettore di 81 interi (0 = vuoto).

(defun prima-cella-vuota (griglia)
  (loop for i from 0 to 80
        when (= (aref griglia i) 0)
        return i))

(defun validi (griglia pos)
  "Valori 1-9 che non violano riga, colonna, box."
  (remove-if (lambda (v) (viola-p griglia pos v))
             '(1 2 3 4 5 6 7 8 9)))

(defun risolvi (griglia)
  (let ((pos (prima-cella-vuota griglia)))
    (if (null pos)
        griglia                           ; nessuna cella vuota: soluzione!
        (some (lambda (v)
                (let ((g2 (copy-seq griglia)))
                  (setf (aref g2 pos) v)
                  (risolvi g2)))
              (validi griglia pos)))))
\end{lstlisting}
\end{esempio}

\begin{intuizione}
\fn{some} restituisce il primo valore non-\fn{nil} restituito
dalla sua funzione. Qui: il primo ramo che porta a una soluzione.
Non esplora gli altri rami una volta trovata la soluzione.
È il modo idiomatico in Lisp per esprimere ``trovami la prima
soluzione che funziona''.
\end{intuizione}

\begin{nota}
Il risolutore di questo esempio è corretto ma non ottimizzato.
Per Sudoku la strategia ``scegli la cella con meno valori validi''
(MRV, minimum remaining values) riduce drasticamente l'albero
di ricerca. La struttura dell'algoritmo non cambia: cambia solo
quale cella viene scelta a ogni passo.

\medskip
La memoization (cache dei risultati già calcolati) non si applica
al backtracking in senso stretto, perché ogni stato è unico.
Si applica invece alla programmazione dinamica, dove sottoproblemi
identici si ripresentano: vedi Appendice~\ref{app:adt} per
una discussione sulle strutture dati adatte.
\end{nota}

\begin{workbook}
\textbf{Esercizi 133--135}\\[4pt]
L'esercizio~133 implementa \fn{viola-p} per il Sudoku.
L'esercizio~134 completa il risolutore.
L'esercizio~135 aggiunge la strategia MRV.\\[4pt]
Questi tre esercizi sono collaterali: non sono nel percorso
minimo, ma chi li fa capisce concretamente come il backtracking
scala su problemi reali.
\end{workbook}

% ================================================================
\section*{Riepilogo del capitolo}
% ================================================================

\begin{tcolorbox}[
  enhanced,
  colback=csBlue!5, colframe=csBlue!40,
  arc=4pt, boxrule=0.6pt,
  left=14pt, right=14pt, top=12pt, bottom=12pt,
  before skip=16pt, after skip=16pt
]
\textbf{Quello che hai imparato:}

\medskip
\begin{itemize}
  \item \textbf{Generate-and-test}: genera tutti i candidati, filtra
        quelli che soddisfano il vincolo. Corretto ma potenzialmente lento.
  \item \textbf{Stato}: la configurazione parziale costruita finora,
        passata come parametro ricorsivo senza mutarla.
  \item \textbf{Scelta}: l'insieme di candidati al passo corrente,
        che dipende dallo stato.
  \item \textbf{Vincolo}: predicato di correttezza, applicabile
        agli stati parziali per il pruning.
  \item \textbf{Pruning}: non espandere i rami che non possono portare
        a soluzioni. Preserva la correttezza, riduce il tempo.
  \item \textbf{Permutazioni e power set}: due istanze canoniche
        del modello, costruite con \fn{mapcan}.
  \item \textbf{N-Queens}: stato = colonne scelte per riga,
        pruning = \fn{sicura-p}. La versione funzionale non modifica
        mai la griglia.
  \item \textbf{\fn{some}}: restituisce la prima soluzione trovata,
        interrompendo la ricerca. Idioma Lisp per ``prima soluzione''.
  \item \textbf{Backtrack implicito}: in Lisp il ritorno dalla
        chiamata ricorsiva \emph{è} il backtrack. Non serve
        stato globale da ripristinare.
\end{itemize}

\medskip
\textbf{Prossimo:} Capitolo~8, \emph{Stringhe, I/O e parsing}.
Dati che vengono dall'esterno: file, terminale, stringhe da interpretare.
\end{tcolorbox}

\begin{workbook}
\textbf{Prima di andare al Capitolo~8:}\\[4pt]
Completa gli esercizi \textbf{121--132} (core del capitolo).
Gli esercizi 133--135 (Sudoku) sono facoltativi ma consigliati.\\[4pt]
Priorità: \textbf{123, 126, 127, 129, 131}.
Questi cinque coprono generate-and-test, stato/scelta/vincolo,
permutazioni, N-Queens e la ricerca della prima soluzione.
\end{workbook}
